% Source: Evans, "Partial Differential Equations" (2nd ed., AMS Graduate Studies in Mathematics vol. 19, 2010),
% Chapter 1 "Introduction", PDF pages 14-26 of the cited source.
% Transcribed from the cited source and organized according to
% leanblueprint conventions (semantic \label, bracketed titles, \uses dependency edges) below.


This chapter surveys the principal theoretical issues concerning the solving of partial differential equations.

To follow the subsequent discussion, the reader should first of all turn to Appendix A and look over the notation presented there, particularly the multiindex notation for partial derivatives.

\section{Partial Differential Equations}
\label{sec:pde-introduction}

A \textit{partial differential equation} (PDE) is an equation involving an unknown function of two or more variables and certain of its partial derivatives.

Using the notation explained in Appendix A, we can write out symbolically a typical PDE, as follows. Fix an integer $k \geq 1$ and let $U$ denote an open subset of $\R^n$.

\begin{definition}[$k$th-Order Partial Differential Equation]
\label{def:kth-order-pde}
\dcref{ch1:1.1-def-1}
\group{PDE Basics}
\lean{EvansLib.IsPDESolutionOn}
\leanok
An expression of the form
\begin{equation}
F(D^k u(x), D^{k-1} u(x), \ldots, Du(x), u(x), x) = 0 \qquad (x \in U)
\tag{1}
\end{equation}
is called a $k$th-order partial differential equation, where
$$F : \R^{n^k} \times \R^{n^{k-1}} \times \cdots \times \R^n \times \R \times U \to \R$$
is given, and
$$u : U \to \R$$
is the unknown.
\end{definition}

We solve the PDE if we find all $u$ verifying (1), possibly only among those functions satisfying certain auxiliary boundary conditions on some part $\Gamma$ of $\partial U$. By finding the solutions we mean, ideally, obtaining simple, \textbf{explicit} solutions, or, failing that, deducing the existence and other properties of solutions.

\begin{definition}[Linear, Semilinear, Quasilinear, and Fully Nonlinear PDE]
\label{def:linear-semilinear-quasilinear-fully-nonlinear}
\dcref{ch1:1.1-def-2}
\group{PDE Basics}
\uses{def:kth-order-pde}
\lean{EvansLib.IsLinearPDE,EvansLib.IsHomogeneousLinearPDE,EvansLib.IsSemilinearPDE,EvansLib.IsQuasilinearPDE,EvansLib.IsFullyNonlinearPDE}
\leanok
(i) The partial differential equation (1) is called \textit{linear} if it has the form
$$\sum_{|\alpha| \leq k} a_\alpha(x) D^\alpha u = f(x)$$
for given functions $a_\alpha$ ($|\alpha| \leq k$), $f$. This linear PDE is \textit{homogeneous} if $f \equiv 0$.

(ii) The PDE (1) is \textit{semilinear} if it has the form
$$\sum_{|\alpha|=k} a_\alpha(x) D^\alpha u + a_0(D^{k-1}u, \ldots, Du, u, x) = 0.$$

(iii) The PDE (1) is \textit{quasilinear} if it has the form
$$\sum_{|\alpha|=k} a_\alpha(D^{k-1}u, \ldots, Du, u, x) D^\alpha u + a_0(D^{k-1}u, \ldots, Du, u, x) = 0.$$

(iv) The PDE (1) is \textit{fully nonlinear} if it depends \textit{nonlinearly} upon the highest order derivatives.
\end{definition}

A \textit{system} of partial differential equations is, informally speaking, a collection of several PDE for several unknown functions.

\begin{definition}[$k$th-Order System of Partial Differential Equations]
\label{def:kth-order-system-pde}
\dcref{ch1:1.1-def-3}
\group{PDE Basics}
\uses{def:kth-order-pde}
\lean{EvansLib.IsPDESystemSolutionOn}
\leanok
An expression of the form
\begin{equation}
\mathbf{F}(D^k \mathbf{u}(x), D^{k-1}\mathbf{u}(x), \ldots, D\mathbf{u}(x), \mathbf{u}(x), x) = 0 \qquad (x \in U)
\tag{2}
\end{equation}
is called a $k^{\text{th}}$-order \textit{system of partial differential equations}, where
$$\mathbf{F} : \R^{mn^k} \times \R^{mn^{k-1}} \times \cdots \times \R^{mn} \times \R^m \times U \to \R^m$$
is given and
$$\mathbf{u} : U \to \R^m, \quad \mathbf{u} = (u^1, \ldots, u^m)$$
is the unknown.
\end{definition}

Here we are supposing that the system comprises the same number $m$ of scalar equations as unknowns $(u^1, \ldots, u^m)$. This is the most common circumstance, although other systems may have fewer or more equations than unknowns.

Systems are classified in the obvious way as being linear, semilinear, etc.

\begin{definition}[Classification of Systems]
\label{def:system-linear-semilinear-quasilinear-fully-nonlinear}
\dcref{ch1:1.1-def-2-systems}
\group{PDE Basics}
\uses{def:kth-order-system-pde,def:linear-semilinear-quasilinear-fully-nonlinear}
\lean{EvansLib.IsLinearPDESystem,EvansLib.IsHomogeneousLinearPDESystem,EvansLib.IsSemilinearPDESystem,EvansLib.IsQuasilinearPDESystem,EvansLib.IsFullyNonlinearPDESystem}
\leanok
A system $\mathbf{F}(D^k\mathbf{u}, \ldots, \mathbf{u}, x) = 0$ valued in $\R^m$ is classified exactly as in the scalar case, with the coefficients now $\R^m$-valued: it is \emph{linear} if $\mathbf{F} = \sum_{|\alpha|\le k} A_\alpha(x) D^\alpha\mathbf{u} - \mathbf{f}(x)$ (matrix coefficients $A_\alpha$), \emph{semilinear} if the top-order coefficients depend only on $x$, \emph{quasilinear} if the top-order coefficients may also depend on the lower-order derivatives, and \emph{fully nonlinear} otherwise. The inclusions homogeneous-linear $\subseteq$ linear $\subseteq$ semilinear $\subseteq$ quasilinear $\subseteq \neg$ fully-nonlinear hold.
\end{definition}

\begin{definition}[Non-Square and Overdetermined Linear Systems]
\label{def:non-square-linear-system}
\group{PDE Basics}
\uses{def:kth-order-system-pde,def:system-linear-semilinear-quasilinear-fully-nonlinear}
\lean{EvansLib.IsLinearPDESystemGen,EvansLib.IsHomogeneousLinearPDESystemGen}
\leanok
As noted above, a system may have more or fewer equations than unknowns. In general a $k$th-order system with unknown $\mathbf{u} : U \to \R^{m}$ whose equations are valued in $\R^{m'}$ (possibly $m' \neq m$) is \emph{linear} if
$$\mathbf{F}(D^k\mathbf{u}, \ldots, \mathbf{u}, x) = \sum_{|\alpha|\le k} A_\alpha(x) D^\alpha\mathbf{u} - \mathbf{f}(x)$$
for rectangular coefficient maps $A_\alpha(x) : \R^{m} \to \R^{m'}$ and source $\mathbf{f} : U \to \R^{m'}$, and \emph{homogeneous} if $\mathbf{f} \equiv 0$. The square case $m' = m$ recovers \cref{def:system-linear-semilinear-quasilinear-fully-nonlinear}; the \emph{overdetermined} case $m' > m$ arises for Maxwell's equations.
\end{definition}

\begin{remark}[``PDE'' as Singular and Plural]
\label{rem:pde-abbreviation}
\dcref{ch1:1.1-rem-1}
\group{PDE Basics}
We use ``PDE'' as an abbreviation for both ``partial differential equation'' and ``partial differential equations''. $\square$
\end{remark}

\section{Examples}
\label{sec:pde-examples}

There is no general theory known concerning the solvability of all partial differential equations. Such a theory is extremely unlikely to exist, given the rich variety of physical, geometric, and probabilistic phenomena which can be modeled by PDE. Instead, research focuses on various particular partial differential equations that are important for applications within and outside of mathematics, with the hope that insight from the origins of these PDE can give clues as to their solutions.

Following is a list of many specific partial differential equations of interest in current research. This listing is intended merely to familiarize the reader with the names and forms of various famous PDE. To display most clearly the mathematical structure of these equations, we have mostly set relevant physical constants to unity. We will later discuss the origin and interpretation of many of these PDE.

Throughout $x \in U$, where $U$ is an open subset of $\R^n$, and $t \geq 0$. Also $Du = D_x u = (u_{x_1}, \ldots, u_{x_n})$ denotes the gradient of $u$ with respect to the spatial variable $x = (x_1, \ldots, x_n)$.

\subsection{Single Partial Differential Equations}
\label{sec:single-pdes}

\subsubsection{Linear Equations}
\label{sec:examples-linear-equations}

\begin{example}[Laplace's Equation]
\label{ex:laplaces-equation}
\dcref{ch1:1.2.1.a.1}
\group{Examples of PDE}
\uses{def:kth-order-pde,def:linear-semilinear-quasilinear-fully-nonlinear}
\lean{EvansLib.laplaceSymbol}
\leanok
$$\Delta u = \sum_{i=1}^n u_{x_i x_i} = 0.$$
\end{example}

\begin{example}[Helmholtz's (or Eigenvalue) Equation]
\label{ex:helmholtz-eigenvalue-equation}
\dcref{ch1:1.2.1.a.2}
\group{Examples of PDE}
\uses{def:kth-order-pde,def:linear-semilinear-quasilinear-fully-nonlinear}
\lean{EvansLib.helmholtzSymbol}
\leanok
$$-\Delta u = \lambda u.$$
\end{example}

\begin{example}[Linear Transport Equation]
\label{ex:linear-transport-equation}
\dcref{ch1:1.2.1.a.3}
\group{Examples of PDE}
\uses{def:kth-order-pde,def:linear-semilinear-quasilinear-fully-nonlinear}
\lean{EvansLib.transportSymbol}
\leanok
$$u_t + \sum_{i=1}^n b^i u_{x_i} = 0.$$
\end{example}

\begin{example}[Liouville's Equation]
\label{ex:liouvilles-equation}
\dcref{ch1:1.2.1.a.4}
\group{Examples of PDE}
\uses{def:kth-order-pde,def:linear-semilinear-quasilinear-fully-nonlinear}
\lean{EvansLib.liouvilleSymbol}
\leanok
$$u_t - \sum_{i=1}^n (b^i u)_{x_i} = 0.$$
\end{example}

\begin{example}[Heat (or Diffusion) Equation]
\label{ex:heat-diffusion-equation}
\dcref{ch1:1.2.1.a.5}
\group{Examples of PDE}
\uses{def:kth-order-pde,def:linear-semilinear-quasilinear-fully-nonlinear}
\lean{EvansLib.heatSymbol}
\leanok
$$u_t - \Delta u = 0.$$
\end{example}

\begin{example}[Schr\"odinger's Equation]
\label{ex:schrodingers-equation}
\dcref{ch1:1.2.1.a.6}
\group{Examples of PDE}
\uses{def:kth-order-pde,def:system-linear-semilinear-quasilinear-fully-nonlinear}
\lean{EvansLib.schrodingerSystemSymbol}
\leanok
$$iu_t + \Delta u = 0.$$
Complex-valued; writing $u = u^1 + i u^2$ splits it into the real linear $2\times 2$ system $\Delta u^1 - u^2_t = 0$, $\Delta u^2 + u^1_t = 0$ (i.e. $\Delta\mathbf{u} + i\,\mathbf{u}_t = 0$, with $i$ acting as the rotation $(a,b)\mapsto(-b,a)$), which is a second-order homogeneous linear system.
\end{example}

\begin{example}[Kolmogorov's Equation]
\label{ex:kolmogorovs-equation}
\dcref{ch1:1.2.1.a.7}
\group{Examples of PDE}
\uses{def:kth-order-pde,def:linear-semilinear-quasilinear-fully-nonlinear}
\lean{EvansLib.kolmogorovSymbol}
\leanok
$$u_t - \sum_{i,j=1}^n a^{ij} u_{x_i x_j} + \sum_{i=1}^n b^i u_{x_i} = 0.$$
\end{example}

\begin{example}[Fokker--Planck Equation]
\label{ex:fokker-planck-equation}
\dcref{ch1:1.2.1.a.8}
\group{Examples of PDE}
\uses{def:kth-order-pde,def:linear-semilinear-quasilinear-fully-nonlinear}
\lean{EvansLib.fokkerPlanckSymbol}
\leanok
$$u_t - \sum_{i,j=1}^n (a^{ij} u)_{x_i x_j} - \sum_{i=1}^n (b^i u)_{x_i} = 0.$$
\end{example}

\begin{example}[Wave Equation]
\label{ex:wave-equation-example}
\dcref{ch1:1.2.1.a.9}
\group{Examples of PDE}
\uses{def:kth-order-pde,def:linear-semilinear-quasilinear-fully-nonlinear}
\lean{EvansLib.waveSymbol}
\leanok
$$u_{tt} - \Delta u = 0.$$
\end{example}

\begin{example}[Telegraph Equation]
\label{ex:telegraph-equation}
\dcref{ch1:1.2.1.a.10}
\group{Examples of PDE}
\uses{def:kth-order-pde,def:linear-semilinear-quasilinear-fully-nonlinear}
\lean{EvansLib.telegraphSymbol}
\leanok
$$u_{tt} + d u_t - u_{xx} = 0.$$
\end{example}

\begin{example}[General Wave Equation]
\label{ex:general-wave-equation}
\dcref{ch1:1.2.1.a.11}
\group{Examples of PDE}
\uses{def:kth-order-pde,def:linear-semilinear-quasilinear-fully-nonlinear}
\lean{EvansLib.generalWaveSymbol}
\leanok
$$u_{tt} - \sum_{i,j=1}^n a^{ij} u_{x_i x_j} + \sum_{i=1}^n b^i u_{x_i} = 0.$$
\end{example}

\begin{example}[Airy's Equation]
\label{ex:airys-equation}
\dcref{ch1:1.2.1.a.12}
\group{Examples of PDE}
\uses{def:kth-order-pde,def:linear-semilinear-quasilinear-fully-nonlinear}
\lean{EvansLib.airySymbol}
\leanok
$$u_t + u_{xxx} = 0.$$
\end{example}

\begin{example}[Beam Equation]
\label{ex:beam-equation}
\dcref{ch1:1.2.1.a.13}
\group{Examples of PDE}
\uses{def:kth-order-pde,def:linear-semilinear-quasilinear-fully-nonlinear}
\lean{EvansLib.beamSymbol}
\leanok
$$u_t + u_{xxxx} = 0.$$
\end{example}

\subsubsection{Nonlinear Equations}
\label{sec:examples-nonlinear-equations}

\begin{example}[Eikonal Equation]
\label{ex:eikonal-equation}
\dcref{ch1:1.2.1.b.1}
\group{Examples of PDE}
\uses{def:kth-order-pde,def:linear-semilinear-quasilinear-fully-nonlinear}
\lean{EvansLib.eikonalSymbol}
\leanok
$$|Du| = 1.$$
\end{example}

\begin{example}[Nonlinear Poisson Equation]
\label{ex:nonlinear-poisson-equation}
\dcref{ch1:1.2.1.b.2}
\group{Examples of PDE}
\uses{def:kth-order-pde,def:linear-semilinear-quasilinear-fully-nonlinear}
\lean{EvansLib.nonlinearPoissonSymbol}
\leanok
$$-\Delta u = f(u).$$
\end{example}

\begin{example}[$p$-Laplacian Equation]
\label{ex:p-laplacian-equation}
\dcref{ch1:1.2.1.b.3}
\group{Examples of PDE}
\uses{def:kth-order-pde,def:linear-semilinear-quasilinear-fully-nonlinear}
\lean{EvansLib.pLaplacianSymbol}
\leanok
$$\Div(|Du|^{p-2}Du) = 0.$$
For a classical solution the divergence expands (chain rule) to
$$|Du|^{p-2}\Delta u + (p-2)|Du|^{p-4} D^2u(Du,Du) = 0,$$
which is quasilinear of order $2$ (the coefficients of $D^2u$ depend on $Du$), and reduces to Laplace's equation at $p = 2$.
\end{example}

\begin{example}[Minimal Surface Equation]
\label{ex:minimal-surface-equation}
\dcref{ch1:1.2.1.b.4}
\group{Examples of PDE}
\uses{def:kth-order-pde,def:linear-semilinear-quasilinear-fully-nonlinear}
\lean{EvansLib.minimalSurfaceSymbol}
\leanok
$$\Div\left(\frac{Du}{(1+|Du|^2)^{1/2}}\right) = 0.$$
For a classical solution the divergence expands to
$$(1+|Du|^2)^{-1/2}\Delta u - (1+|Du|^2)^{-3/2} D^2u(Du,Du) = 0,$$
which is quasilinear of order $2$.
\end{example}

\begin{example}[Monge--Amp\`ere Equation]
\label{ex:monge-ampere-equation}
\dcref{ch1:1.2.1.b.5}
\group{Examples of PDE}
\uses{def:kth-order-pde,def:linear-semilinear-quasilinear-fully-nonlinear}
\lean{EvansLib.mongeAmpereSymbol}
\leanok
$$\det(D^2u) = f.$$
\end{example}

\begin{example}[Hamilton--Jacobi Equation]
\label{ex:hamilton-jacobi-equation}
\dcref{ch1:1.2.1.b.6}
\group{Examples of PDE}
\uses{def:kth-order-pde,def:linear-semilinear-quasilinear-fully-nonlinear}
\lean{EvansLib.hamiltonJacobiSymbol}
\leanok
$$u_t + H(Du, x) = 0.$$
First order. For a nonlinear Hamiltonian $H$ the equation is fully nonlinear: for the archetypal quadratic Hamiltonian $H(p,x) = |p|^2$ (so $u_t + |Du|^2 = 0$) the symbol is quadratic, not affine, in the top-order slot $Du$ (formalized as \lean{EvansLib.hamiltonJacobiQuadraticSymbol_isFullyNonlinearPDE}).
\end{example}

\begin{example}[Scalar Conservation Law]
\label{ex:scalar-conservation-law}
\dcref{ch1:1.2.1.b.7}
\group{Examples of PDE}
\uses{def:kth-order-pde,def:linear-semilinear-quasilinear-fully-nonlinear}
\lean{EvansLib.conservationLawSymbol}
\leanok
$$u_t + \Div\mathbf{F}(u) = 0.$$
For a classical solution the chain rule gives $u_t + \sum_i (F^i)'(u) u_{x_i} = 0$, quasilinear of order $1$ (the top-order coefficients $(F^i)'(u)$ depend on $u$).
\end{example}

\begin{example}[Inviscid Burgers' Equation]
\label{ex:inviscid-burgers-equation}
\dcref{ch1:1.2.1.b.8}
\group{Examples of PDE}
\uses{def:kth-order-pde,def:linear-semilinear-quasilinear-fully-nonlinear,ex:scalar-conservation-law}
\lean{EvansLib.burgersSymbol}
\leanok
$$u_t + uu_x = 0.$$
\end{example}

\begin{example}[Scalar Reaction-Diffusion Equation]
\label{ex:scalar-reaction-diffusion-equation}
\dcref{ch1:1.2.1.b.9}
\group{Examples of PDE}
\uses{def:kth-order-pde,def:linear-semilinear-quasilinear-fully-nonlinear}
\lean{EvansLib.reactionDiffusionSymbol}
\leanok
$$u_t - \Delta u = f(u).$$
\end{example}

\begin{example}[Porous Medium Equation]
\label{ex:porous-medium-equation}
\dcref{ch1:1.2.1.b.10}
\group{Examples of PDE}
\uses{def:kth-order-pde,def:linear-semilinear-quasilinear-fully-nonlinear}
\lean{EvansLib.nonlinearDiffusionSymbol}
\leanok
$$u_t - \Delta(u^\gamma) = 0.$$
For a classical solution the chain rule gives $u_t - \gamma u^{\gamma-1}\Delta u - \gamma(\gamma-1)u^{\gamma-2}|Du|^2 = 0$, quasilinear of order $2$; formalized in the general nonlinear-diffusion form $u_t - \varphi'(u)\Delta u - \varphi''(u)|Du|^2 = 0$ with $\varphi(u) = u^\gamma$.
\end{example}

\begin{example}[Nonlinear Wave Equations]
\label{ex:nonlinear-wave-equations}
\dcref{ch1:1.2.1.b.11}
\group{Examples of PDE}
\uses{def:kth-order-pde,def:linear-semilinear-quasilinear-fully-nonlinear}
\lean{EvansLib.nonlinearWaveSymbol,EvansLib.nonlinearWaveDivSymbol}
\leanok
$$u_{tt} - \Delta u = f(u),$$
$$u_{tt} - \Div\,\mathbf{a}(Du) = 0.$$
The first (reaction) form is semilinear of order $2$; the second (divergence) form expands to $u_{tt} - \sum_{i,j}(\partial a_j/\partial p_i)(Du)\,u_{x_ix_j} = 0$, quasilinear of order $2$.
\end{example}

\begin{example}[Korteweg--de Vries (KdV) Equation]
\label{ex:kdv-equation}
\dcref{ch1:1.2.1.b.12}
\group{Examples of PDE}
\uses{def:kth-order-pde,def:linear-semilinear-quasilinear-fully-nonlinear}
\lean{EvansLib.kdvSymbol}
\leanok
$$u_t + uu_x + u_{xxx} = 0.$$
\end{example}

\subsection{Systems of Partial Differential Equations}
\label{sec:systems-of-pdes}

\subsubsection{Linear Systems}
\label{sec:examples-linear-systems}

\begin{example}[Equilibrium Equations of Linear Elasticity]
\label{ex:linear-elasticity-equilibrium}
\dcref{ch1:1.2.2.a.1}
\group{Examples of PDE}
\uses{def:kth-order-system-pde,def:system-linear-semilinear-quasilinear-fully-nonlinear}
\lean{EvansLib.linearElasticityEquilibriumSymbol}
\leanok
$$\mu\Delta\mathbf{u} + (\lambda + \mu)D(\Div \mathbf{u}) = 0.$$
A second-order linear system for the displacement $\mathbf{u} : \R^n \to \R^n$.
\end{example}

\begin{example}[Evolution Equations of Linear Elasticity]
\label{ex:linear-elasticity-evolution}
\dcref{ch1:1.2.2.a.2}
\group{Examples of PDE}
\uses{def:kth-order-system-pde,def:system-linear-semilinear-quasilinear-fully-nonlinear,ex:linear-elasticity-equilibrium}
\lean{EvansLib.linearElasticityEvolutionSymbol}
\leanok
$$\mathbf{u}_t - \mu\Delta\mathbf{u} - (\lambda + \mu)D(\Div \mathbf{u}) = 0.$$
A second-order linear system: the time derivative $\mathbf{u}_t$ plus the (negated) elastic operator, both linear.
\end{example}

\begin{example}[Maxwell's Equations]
\label{ex:maxwells-equations}
\dcref{ch1:1.2.2.a.3}
\group{Examples of PDE}
\uses{def:kth-order-system-pde,def:non-square-linear-system}
\lean{EvansLib.maxwellSymbol,EvansLib.maxwellSymbol_isLinearPDESystemGen}
\leanok
$$\begin{cases}
\mathbf{E}_t = \curl \mathbf{B} \\
\mathbf{B}_t = -\curl \mathbf{E} \\
\Div \mathbf{B} = \Div \mathbf{E} = 0.
\end{cases}$$
For the electromagnetic field $\mathbf{u} = (\mathbf{E}, \mathbf{B}) : \R^4 \to \R^6$ (with $\mathbf{E} = (u^0,u^1,u^2)$, $\mathbf{B} = (u^3,u^4,u^5)$) this is a first-order homogeneous linear system. It is \emph{overdetermined}: the two curl equations contribute $3+3$ scalar equations and the two divergence constraints $1+1$ more, so there are $8$ equations for the $6$ unknowns; its symbol is valued in $\R^8$.
\end{example}

\subsubsection{Nonlinear Systems}
\label{sec:examples-nonlinear-systems}

\begin{example}[System of Conservation Laws]
\label{ex:system-conservation-laws}
\dcref{ch1:1.2.2.b.1}
\group{Examples of PDE}
\uses{def:kth-order-system-pde,def:system-linear-semilinear-quasilinear-fully-nonlinear,ex:scalar-conservation-law}
\lean{EvansLib.conservationLawSystemSymbol}
\leanok
$$\mathbf{u}_t + \Div \mathbf{F}(\mathbf{u}) = 0.$$
For a classical solution the chain rule gives $\mathbf{u}_t + \sum_i (D\mathbf{F}^i)(\mathbf{u})\,\mathbf{u}_{x_i} = 0$, quasilinear of order $1$ (the top-order coefficient matrices $D\mathbf{F}^i(\mathbf{u})$ depend on $\mathbf{u}$).
\end{example}

\begin{example}[Reaction-Diffusion System]
\label{ex:reaction-diffusion-system}
\dcref{ch1:1.2.2.b.2}
\group{Examples of PDE}
\uses{def:kth-order-system-pde,def:system-linear-semilinear-quasilinear-fully-nonlinear,ex:scalar-reaction-diffusion-equation}
\lean{EvansLib.reactionDiffusionSystemSymbol}
\leanok
$$\mathbf{u}_t - \Delta\mathbf{u} = \mathbf{f}(\mathbf{u}).$$
A second-order semilinear system: the heat operator $\mathbf{u}_t - \Delta\mathbf{u}$ is linear and the reaction $\mathbf{f}(\mathbf{u})$ is lower order.
\end{example}

\begin{example}[Euler's Equations for Incompressible, Inviscid Flow]
\label{ex:euler-equations-incompressible-inviscid}
\dcref{ch1:1.2.2.b.3}
\group{Examples of PDE}
\uses{def:kth-order-system-pde,def:system-linear-semilinear-quasilinear-fully-nonlinear}
\lean{EvansLib.eulerSymbol,EvansLib.eulerSymbol_isQuasilinearPDESystem}
\leanok
$$\begin{cases}
\mathbf{u}_t + \mathbf{u} \cdot D\mathbf{u} = -D p \\
\Div \mathbf{u} = 0.
\end{cases}$$
The velocity--pressure pair $(\mathbf{u}, p) : \R^{n+1} \to \R^{n+1}$ solves an $(n+1) \times (n+1)$ first-order system: the $n$ momentum equations $u^i_t + \sum_j u^j u^i_{x_j} + p_{x_i} = 0$ and incompressibility $\Div \mathbf{u} = \sum_j u^j_{x_j} = 0$. It is quasilinear of order $1$: the convective term carries the top-order derivative $u^i_{x_j}$ with coefficient $u^j$ depending on the lower-order jet $\mathbf{u}$.
\end{example}

\begin{example}[Navier--Stokes Equations for Incompressible, Viscous Flow]
\label{ex:navier-stokes-equations-incompressible-viscous}
\dcref{ch1:1.2.2.b.4}
\group{Examples of PDE}
\uses{def:kth-order-system-pde,def:system-linear-semilinear-quasilinear-fully-nonlinear,ex:euler-equations-incompressible-inviscid}
\lean{EvansLib.nsSymbol,EvansLib.nsSymbol_isSemilinearPDESystem}
\leanok
$$\begin{cases}
\mathbf{u}_t + \mathbf{u} \cdot D\mathbf{u} - \Delta\mathbf{u} = -D p \\
\Div \mathbf{u} = 0.
\end{cases}$$
Adding the viscous term $-\Delta\mathbf{u}$ to Euler's equations raises the order to $2$. The top-order operator $-\Delta\mathbf{u}$ is linear with constant coefficients, while the convective nonlinearity $\mathbf{u}\cdot D\mathbf{u}$ is of lower (first) order, so the system is \emph{semilinear} of order $2$; its lower-order part is exactly Euler's operator.
\end{example}

See Zwillinger \cite{Zwillinger1984} for a much more extensive listing of interesting PDE.

\section{Strategies for Studying PDE}
\label{sec:strategies-studying-pde}

As explained in \cref{sec:pde-introduction} our goal is the discovery of ways to solve partial differential equations of various sorts, but---as should now be clear in view of the many diverse examples set forth in \cref{sec:pde-examples}---this is no easy task. And indeed the very question of what it means to ``solve'' a given PDE can be subtle, depending in large part on the particular structure of the problem at hand.

\subsection{Well-Posed Problems, Classical Solutions}
\label{sec:well-posed-classical-solutions}

The informal notion of a well-posed problem captures many of the desirable features of what it means to solve a PDE.

\begin{definition}[Well-Posed Problem]
\label{def:well-posed-problem}
\dcref{ch1:1.3.1-def-1}
\group{Well-Posedness}
\lean{EvansLib.PDEProblem.IsWellPosed}
\leanok
We say that a given problem for a partial differential equation is \textit{well-posed} if

(a) the problem in fact has a solution;

(b) this solution is unique;

and

(c) the solution depends continuously on the data given in the problem.
\end{definition}

The last condition is particularly important for problems arising from physical applications: we would prefer that our (unique) solution changes only a little when the conditions specifying the problem change a little. (For many problems, on the other hand, uniqueness is not to be expected. In these cases the primary mathematical tasks are to classify and characterize the solutions.)

Now clearly it would be desirable to ``solve'' PDE in such a way that (a)--(c) hold. But notice that we still have not carefully defined what we mean by a ``solution''.

\begin{definition}[Classical Solution]
\label{def:classical-solution}
\dcref{ch1:1.3.1-def-2}
\group{Well-Posedness}
\uses{def:kth-order-pde}
\lean{EvansLib.IsClassicalSolutionOn}
\leanok
Should we ask, for example, that a ``solution'' $u$ must be real analytic or at least infinitely differentiable? This might be desirable, but perhaps we are asking too much. Maybe it would be wiser to require a solution of a PDE of order $k$ to be at least $k$ times continuously differentiable. Then at least all the derivatives which appear in the statement of the PDE will exist and be continuous, although maybe certain higher derivatives will not exist. Let us informally call a solution with this much smoothness a \textit{classical} solution of the PDE: this is certainly the most obvious notion of solution.
\end{definition}

So by solving a partial differential equation in the classical sense we mean if possible to write down a formula for a classical solution satisfying (a)--(c) of \cref{def:well-posed-problem} above, or at least to show such a solution exists, and to deduce various of its properties.

\subsection{Weak Solutions and Regularity}
\label{sec:weak-solutions-regularity}

But can we achieve this? The answer is that certain specific partial differential equations (e.g. Laplace's equation) can be solved in the classical sense, but many others, if not most others, cannot. Consider for instance the scalar conservation law
$$u_t + F(u)_x = 0.$$
We will see in \S 3.4 that this PDE governs various one-dimensional phenomena involving fluid dynamics, and in particular models the formation and propagation of shock waves. Now a shock wave is a curve of discontinuity of the solution $u$; and so if we wish to study conservation laws, and recover the underlying physics, we must surely allow for solutions $u$ which are not continuously differentiable or even continuous. In general, as we shall see, the conservation law has no classical solutions, but \textit{is} well-posed if we allow for properly defined \textit{generalized} or \textit{weak} solutions.

This is all to say that we may be forced by the structure of the particular equation to abandon the search for smooth, classical solutions. We must instead, while still hoping to achieve the well-posedness conditions (a)--(c) of \cref{def:well-posed-problem}, investigate a wider class of candidates for solutions. And in fact, even for those PDE which turn out to be classically solvable, it is often most expedient initially to search for some appropriate kind of weak solution.

The point is this: if from the outset we demand that our solutions be very regular, say $k$-times continuously differentiable, then we are usually going to have a really hard time finding them, as our proofs must then necessarily include possibly intricate demonstrations that the functions we are building are in fact smooth enough. A far more reasonable strategy is to consider as separate the \textit{existence} and the \textit{smoothness} (or \textit{regularity}) problems. The idea is to define for a given PDE a reasonably wide notion of a \textit{weak solution}, with the expectation that since we are not asking too much by way of smoothness of this weak solution, it may be easier to establish its existence, uniqueness, and continuous dependence on the given data. Thus, to repeat, it is often wise to aim at proving well-posedness in some appropriate class of weak or generalized solutions.

Now, as noted above, for various partial differential equations this is the best that can be done. For other equations we can hope that our weak solution may turn out after all to be smooth enough to qualify as a classical solution. This leads to the question of \textit{regularity} of weak solutions. As we will see, it is often the case that the existence of weak solutions depends upon rather simple estimates plus ideas of functional analysis, whereas the regularity of the weak solutions, when true, usually rests upon many intricate calculus estimates.

Let me explicitly note here that once we are past Part I (Chapters 2--4), our efforts will be largely devoted to proving mathematically the existence of solutions to various sorts of partial differential equations, and not so much to deriving formulas for these solutions. This may seem wasted or misguided effort, but in fact mathematicians are like theologians: we regard existence as the prime attribute of what we study. But unlike most theologians, we need not always rely upon faith alone.

\subsection{Typical Difficulties}
\label{sec:typical-difficulties}

Following are some vague but general principles, which may be useful to keep in mind:
\begin{enumerate}
\item Nonlinear equations are more difficult than linear equations; and, indeed, the more the nonlinearity affects the higher derivatives, the more difficult the PDE is.
\item Higher-order PDE are more difficult than lower-order PDE.
\item Systems are harder than single equations.
\item Partial differential equations entailing many independent variables are harder than PDE entailing few independent variables.
\item For most partial differential equations it is not possible to write out explicit formulas for solutions.
\end{enumerate}

None of these assertions is without important exceptions.

\section{Overview}
\label{sec:chapter1-overview}

This textbook is divided into three major Parts.

\textbf{PART I: Representation Formulas for Solutions}

Here we identify those important partial differential equations for which in certain circumstances explicit or more-or-less explicit formulas can be had for solutions. The general progression of the exposition is from direct formulas for certain linear equations, to far less concrete representation formulas, of a sort, for various nonlinear PDE.

Chapter 2 is a detailed study of four exactly solvable partial differential equations: the linear transport equation, Laplace's equation, the heat equation, and the wave equation. These PDE, which serve as archetypes for the more complicated equations introduced later, admit directly computable solutions, at least in the case that there is no domain whose boundary geometry complicates matters. The explicit formulas are augmented by various indirect, but easy and attractive, ``energy''-type arguments, which serve as motivation for the developments in Chapters 6, 7 and thereafter.

Chapter 3 continues the theme of searching for explicit formulas, now for general first-order nonlinear PDE. The key insight is that such PDE can, locally at least, be transformed into systems of ordinary differential equations (ODE), the characteristic equations. We stipulate that once the problem becomes ``only'' the question of integrating a system of ODE, it is in principle solved, \textit{sometimes quite explicitly}. The derivation of the characteristic equations given in the text is very simple and does not require any geometric insights. It is in truth so easy to derive the characteristic equations that no real purpose is had by dealing with the quasilinear case first.

We introduce also the Hopf--Lax formula for Hamilton--Jacobi equations (\S 3.3) and the Lax--Oleinik formula for scalar conservation laws (\S 3.4). (Some knowledge of measure theory is useful here, but is not essential.) These sections provide an early acquaintance with the global theory of these important nonlinear PDE, and so motivate the later Chapters 10 and 11.

Chapter 4 is a grab bag of techniques for explicitly (or kind of explicitly) solving various linear and nonlinear partial differential equations, and the reader should study only whatever seems interesting. The section on the Fourier transform is, however, essential. The Cauchy--Kovalevskaya Theorem appears at the very end. Although this is basically the only general existence theorem in the subject, and thus logically should perhaps be regarded as central, in practice these power series methods are not so prevalent.

\textbf{PART II: Theory for Linear Partial Differential Equations}

Next we abandon the search for explicit formulas and instead rely on functional analysis and relatively easy ``energy'' estimates to prove the existence of weak solutions to various linear PDE. We investigate also the uniqueness and regularity of such solutions, and deduce various other properties.

Chapter 5 is an introduction to Sobolev spaces, the proper setting for the study of many linear and nonlinear partial differential equations via energy methods. This is a hard chapter, the real worth of which is only later revealed, and requires some basic knowledge of Lebesgue measure theory. However, the requirements are not really so great, and the review in Appendix E should suffice. In my opinion there is no particular advantage in considering only the Sobolev spaces with exponent $p = 2$, and indeed insisting upon this obscures the two central inequalities, those of Gagliardo--Nirenberg--Sobolev (\S 5.6.1) and of Morrey (\S 5.6.2).

In Chapter 6 we vastly generalize our knowledge of Laplace's equation to other second-order elliptic equations. Here we work through a rather complete treatment of existence, uniqueness and regularity theory for solutions, including the maximum principle, and also a reasonable introduction to the study of eigenvalues, including a discussion of the principal eigenvalue for nonselfadjoint operators.

Chapter 7 expands the energy methods to a variety of linear partial differential equations characterizing evolutions in time. We broaden our earlier investigation of the heat equation to general second-order parabolic PDE, and of the wave equation to general second-order hyperbolic PDE. We study as well linear first-order hyperbolic systems, with the aim of motivating the developments concerning nonlinear systems of conservation laws in Chapter 11. The concluding section 7.4 presents the alternative functional analytic method of semigroups for building solutions.

(Missing from this long Part II on linear partial differential equations is any discussion of distribution theory or potential theory. These are important topics, but for our purposes seem dispensable, even in a book of such length. These omissions do not slow us up much, and make room for more nonlinear theory.)

\textbf{PART III: Theory for Nonlinear Partial Differential Equations}

This section parallels for nonlinear PDE the development in Part II, but is far less unified in its approach, as the various types of nonlinearity must be treated in quite different ways.

Chapter 8 commences the general study of nonlinear partial differential equations with an extensive discussion of the calculus of variations. Here we set forth a careful derivation of the direct method for deducing the existence of minimizers, and discuss also a variety of variational systems and constrained problems, as well as minimax methods. Variational theory is the most useful and accessible of the methods for nonlinear PDE, and so this chapter is fundamental.

Chapter 9 is, rather like Chapter 4 before, a gathering of assorted other techniques of use for nonlinear elliptic and parabolic partial differential equations. We encounter here monotonicity and fixed-point methods, and a variety of other devices, mostly involving the maximum principle. We study as well certain nice aspects of nonlinear semigroup theory, to complement the linear semigroup theory from Chapter 7.

Chapter 10 is an introduction to the modern theory of Hamilton-Jacobi PDE, and in particular to the notion of ``viscosity solutions''. We encounter also the connections with the optimal control of ODE, through dynamic programming.

Chapter 11 picks up from Chapter 3 the discussion of conservation laws, now systems of conservation laws. Unlike the general theoretical developments in Chapters 5--9, for which Sobolev spaces provide the proper abstract framework, we are forced to employ here direct linear algebra and calculus computations. We pay particular attention to the solution of Riemann's problem and to entropy criteria.

Appendices A--E provide for the reader's convenience some background material, with selected proofs, on inequalities, linear functional analysis, measure theory, etc.

The Bibliography primarily provides a listing of interesting PDE books to consult for further information. Since this is a textbook, and not a reference monograph, I have mostly not attempted to track down and document the original sources for the myriads of ideas and methods we will encounter. The mathematical literature for partial differential equations is truly vast, but the books cited in the Bibliography should at least provide a starting point for locating the primary sources.

\section{Problems}
\label{sec:chapter1-problems}

%ORIGINAL-NUMBER: 1
\begin{exercise}
Classify each of the partial differential equations in \cref{sec:pde-examples} as follows:
\begin{enumerate}
\item[(a)] Is the PDE linear, semilinear, quasilinear or fully nonlinear?
\item[(b)] What is the order of the PDE?
\end{enumerate}
\end{exercise}

The next exercises provide some practice with the multiindex notation introduced in Appendix A.

%ORIGINAL-NUMBER: 2
\begin{proposition}[Multinomial Theorem]
\label{ex:multinomial-theorem}
\group{Multiindex Calculus}
\lean{EvansLib.multinomial_theorem}
\leanok
For $x = (x_1, \ldots, x_n)$ and $k \in \N$, the \textit{Multinomial Theorem} states
$$\left(x_1 + \ldots + x_n\right)^k = \sum_{|\alpha|=k} \binom{|\alpha|}{\alpha} x^\alpha,$$
where $\binom{|\alpha|}{\alpha} := \frac{|\alpha|!}{\alpha!}$, $\alpha! = \alpha_1!\alpha_2! \cdots \alpha_n!$, and $x^\alpha = x_1^{\alpha_1} \cdots x_n^{\alpha_n}$. The sum is taken over all multiindices $\alpha = (\alpha_1, \ldots, \alpha_n)$ with $|\alpha| = k$.
\end{proposition}
\begin{proof}
\leanok
Expand the $k$-fold product $(x_1 + \cdots + x_n)^k = \prod_{r=1}^{k}\bigl(\sum_{i=1}^n x_i\bigr)$ by the distributive law: it equals $\sum_{f} \prod_{r=1}^{k} x_{f(r)}$, summed over all functions $f : \{1,\ldots,k\} \to \{1,\ldots,n\}$. For such an $f$ put $\alpha_i := \#\{r : f(r) = i\}$; then $\alpha = (\alpha_1,\ldots,\alpha_n)$ is a multiindex with $|\alpha| = \sum_i \alpha_i = k$ and $\prod_{r} x_{f(r)} = x^\alpha$. Conversely, the number of functions $f$ giving rise to a fixed multiindex $\alpha$ is the number of ways to split the $k$ positions into ordered groups of sizes $\alpha_1, \ldots, \alpha_n$, namely $\frac{k!}{\alpha_1!\cdots\alpha_n!} = \binom{|\alpha|}{\alpha}$. Grouping the sum over $f$ by its associated $\alpha$ therefore gives $\sum_{|\alpha|=k} \binom{|\alpha|}{\alpha} x^\alpha$.
\end{proof}

\begin{definition}[Partial Derivative]
\label{def:partial-derivative}
\group{Multiindex Calculus}
\lean{EvansLib.partialDeriv}
\leanok
For $f : \R^n \to \R$ and an index $i \in \{1, \ldots, n\}$, the \textit{$i$th
partial derivative} is the directional derivative of $f$ along the $i$th standard
basis vector $e_i$:
$$\partial_{x_i} f (x) = Df(x)(e_i).$$
\end{definition}

\begin{definition}[Multiindex Partial Derivative]
\label{def:multiindex-partial-derivative}
\group{Multiindex Calculus}
\uses{def:partial-derivative}
\lean{EvansLib.multiPartial}
\leanok
For a multiindex $\alpha = (\alpha_1, \ldots, \alpha_n) \in \N^n$ and a function
$f : \R^n \to \R$, the \textit{multiindex partial derivative} is
$$D^\alpha f = \frac{\partial^{|\alpha|} f}{\partial x_1^{\alpha_1} \cdots \partial x_n^{\alpha_n}}
  = \partial_{x_1}^{\alpha_1} \cdots \partial_{x_n}^{\alpha_n} f,$$
each coordinate partial $\partial_{x_i}$ being applied $\alpha_i$ times, and
$|\alpha| = \alpha_1 + \cdots + \alpha_n$.
\end{definition}

\begin{proposition}[Equality of Mixed Partial Derivatives]
\label{prop:clairaut-partial-commute}
\group{Multiindex Calculus}
\uses{def:partial-derivative,def:multiindex-partial-derivative}
\lean{EvansLib.partialDeriv_comm}
\leanok
For smooth $f : \R^n \to \R$ and indices $i, j$,
$$\partial_{x_i}\partial_{x_j} f = \partial_{x_j}\partial_{x_i} f.$$
Consequently $D^\alpha f$ does not depend on the order in which the coordinate
partials are applied.
\end{proposition}
\begin{proof}
\leanok
Since $\partial_{x_j} f = Df(\cdot)(e_j)$, differentiating again gives, for the
$C^2$ function $f$,
$$\partial_{x_i}\partial_{x_j} f(x) = D\big(Df(\cdot)(e_j)\big)(x)(e_i)
   = D^2 f(x)(e_i, e_j),$$
where $D^2 f(x)$ is the second Fréchet derivative, a bilinear form. By the
Clairaut--Schwarz theorem $D^2 f(x)$ is symmetric, so
$D^2 f(x)(e_i, e_j) = D^2 f(x)(e_j, e_i) = \partial_{x_j}\partial_{x_i} f(x)$.
Any two orderings of the coordinate partials in $D^\alpha f$ differ by a sequence
of adjacent transpositions, each of which is an instance of this equality; hence
$D^\alpha f$ is well defined.
\end{proof}

\begin{proposition}[Iterated Leibniz Rule Along One Axis]
\label{prop:iterated-leibniz-one-axis}
\group{Multiindex Calculus}
\uses{def:partial-derivative}
\lean{EvansLib.partialDeriv_iterate_mul}
\leanok
For smooth $u, v : \R^n \to \R$, an axis $i$ and $m \in \N$,
$$\partial_{x_i}^m(uv) = \sum_{j=0}^m \binom{m}{j}\,
  (\partial_{x_i}^j u)\,(\partial_{x_i}^{m-j} v).$$
\end{proposition}
\begin{proof}
\leanok
Fix $x$ and set $\gamma(t) = x + t e_i$. The one-variable restrictions
$\tilde u(t) = u(\gamma(t))$ and $\tilde v(t) = v(\gamma(t))$ are smooth, and by
the chain rule $\tilde u^{(k)}(t) = (\partial_{x_i}^k u)(\gamma(t))$ for every $k$
(and likewise for $v$ and for the product $uv$, since restriction to the line
commutes with multiplication). The one-variable Leibniz rule
$(\tilde u\,\tilde v)^{(m)} = \sum_{j=0}^m \binom{m}{j}\tilde u^{(j)}\tilde v^{(m-j)}$,
evaluated at $t = 0$, is exactly the claimed identity at $x$.
\end{proof}

%ORIGINAL-NUMBER: 3
\begin{proposition}[Leibniz' Formula]
\label{prop:leibniz-formula}
\group{Multiindex Calculus}
\uses{def:multiindex-partial-derivative,prop:clairaut-partial-commute,prop:iterated-leibniz-one-axis}
\lean{EvansLib.multiPartial_mul}
\leanok
For smooth $u,v : \R^n \to \R$,
$$D^\alpha(uv) = \sum_{\beta \leq \alpha} \binom{\alpha}{\beta} D^\beta u \, D^{\alpha-\beta} v,$$
where $\binom{\alpha}{\beta} := \frac{\alpha!}{\beta!(\alpha-\beta)!} = \prod_{i=1}^n \binom{\alpha_i}{\beta_i}$, and $\beta \leq \alpha$ means $\beta_i \leq \alpha_i$ $(i = 1, \ldots, n)$.
\end{proposition}
\begin{proof}
\leanok
Induct on $|\alpha|$. If $|\alpha| = 0$ both sides equal $uv$. Otherwise pick $i$
with $\alpha_i \geq 1$ and write $D^\alpha = \partial_{x_i} D^{\alpha - e_i}$,
which is legitimate by \cref{prop:clairaut-partial-commute}. By the inductive
hypothesis,
$$D^{\alpha - e_i}(uv) = \sum_{\gamma \leq \alpha - e_i}
   \binom{\alpha - e_i}{\gamma} (D^\gamma u)(D^{\alpha - e_i - \gamma} v).$$
Applying $\partial_{x_i}$ and the first-order product rule
$\partial_{x_i}(gh) = (\partial_{x_i}g)h + g(\partial_{x_i}h)$ to each summand,
then using $\partial_{x_i} D^\gamma = D^{\gamma + e_i}$, the coefficient of
$(D^\beta u)(D^{\alpha - \beta} v)$ becomes
$\binom{\alpha - e_i}{\beta - e_i} + \binom{\alpha - e_i}{\beta}$. As
$\binom{\alpha}{\beta} = \prod_k \binom{\alpha_k}{\beta_k}$ and only the $i$th
factor differs, Pascal's rule
$\binom{\alpha_i}{\beta_i} = \binom{\alpha_i - 1}{\beta_i - 1} + \binom{\alpha_i - 1}{\beta_i}$
gives $\binom{\alpha}{\beta}$, completing the induction.
\end{proof}

%ORIGINAL-NUMBER: 4
\begin{proposition}[Directional Derivatives via the Multinomial Theorem]
\label{prop:directional-derivative-multinomial}
\group{Multiindex Calculus}
\uses{def:multiindex-partial-derivative,prop:clairaut-partial-commute,ex:multinomial-theorem}
\lean{EvansLib.iteratedDeriv_scaled_line_multinomial}
\leanok
Let $f : \R^n \to \R$ be smooth and fix $x \in \R^n$. The one-variable function
$g(t) = f(tx)$ satisfies, for every $j \in \N$,
$$g^{(j)}(0) = \sum_{|\alpha|=j} \binom{j}{\alpha} D^\alpha f(0)\, x^\alpha,$$
where $\binom{j}{\alpha} = j!/\alpha!$ is the multinomial coefficient and
$x^\alpha = x_1^{\alpha_1}\cdots x_n^{\alpha_n}$.
\end{proposition}
\begin{proof}
\leanok
Write $D_x h(y) := Dh(y)(x) = \sum_{i=1}^n x_i\,\partial_{x_i} h(y)$ for the
directional derivative of $h$ along $x$. Applying the chain rule to the line
$t \mapsto tx$ shows that the $j$-th ordinary derivative of $g$ is the $j$-fold
directional derivative evaluated along the line, $g^{(j)}(t) = (D_x^{\,j} f)(tx)$;
in particular $g^{(j)}(0) = (D_x^{\,j}f)(0)$. Since the coordinate partials commute
(\cref{prop:clairaut-partial-commute}), the operator power
$D_x^{\,j} = \bigl(\sum_i x_i\partial_{x_i}\bigr)^j$ expands by the Multinomial
Theorem as $D_x^{\,j} = \sum_{|\alpha|=j}\binom{j}{\alpha} x^\alpha D^\alpha$;
evaluating at $0$ gives the claim. The operator expansion is proved by induction
on $j$: differentiating $\sum_{|\alpha|=j}\binom{j}{\alpha}x^\alpha D^\alpha f$ once
more along $x$ sends each $\alpha$ to $\alpha + e_i$ with weight $x_i$, and the
Pascal relation $\sum_{i\,:\,\beta_i\geq 1}\binom{j}{\beta-e_i} = \binom{j+1}{\beta}$
recombines the coefficients.
\end{proof}

\begin{proposition}[Taylor's Formula in Multiindex Notation]
\label{prop:taylor-multiindex}
\group{Multiindex Calculus}
\uses{def:multiindex-partial-derivative,prop:directional-derivative-multinomial}
\lean{EvansLib.taylor_multiindex}
\leanok
Assume that $f : \R^n \to \R$ is smooth. Then for each $k = 1, 2, \ldots$,
$$f(x) = \sum_{|\alpha| \leq k} \frac{1}{\alpha!} D^\alpha f(0) x^\alpha + O(|x|^{k+1}) \quad \text{as } x \to 0.$$
This is \textit{Taylor's formula} in multiindex notation.
\end{proposition}
\begin{proof}
\leanok
Fix $x$ and set $g(t) = f(tx)$, so that $g(1) = f(x)$ and $g$ is smooth. By the
one-variable Taylor theorem with Lagrange remainder there is $\xi \in (0,1)$ with
$$g(1) = \sum_{j=0}^{k}\frac{g^{(j)}(0)}{j!} + \frac{g^{(k+1)}(\xi)}{(k+1)!}.$$
By \cref{prop:directional-derivative-multinomial} and
$\binom{j}{\alpha}/j! = 1/\alpha!$ (for $|\alpha| = j$),
$$\frac{g^{(j)}(0)}{j!} = \sum_{|\alpha|=j}\frac{1}{\alpha!}D^\alpha f(0)\, x^\alpha,$$
so summing over $j \leq k$ produces the multiindex Taylor polynomial. For the
remainder, \cref{prop:directional-derivative-multinomial} at the base point
$\xi x$ gives
$g^{(k+1)}(\xi) = \sum_{|\alpha|=k+1}\binom{k+1}{\alpha} D^\alpha f(\xi x)\, x^\alpha$.
Each $D^\alpha f$ is continuous, hence bounded by some constant $M$ on the closed
unit ball; for $|x| \leq 1$ and $\xi \in (0,1)$ the point $\xi x$ lies in that
ball, while $|x^\alpha| \leq |x|^{k+1}$ whenever $|\alpha| = k+1$. Hence
$|g^{(k+1)}(\xi)| \leq \bigl(\sum_{|\alpha|=k+1}\binom{k+1}{\alpha} M\bigr)|x|^{k+1}$,
and the remainder is $O(|x|^{k+1})$ as $x \to 0$.
\end{proof}
